\documentclass[12pt]{article}
\usepackage[utf8]{inputenc}
\usepackage[T1]{fontenc}
\usepackage[margin=1in]{geometry}
\usepackage{setspace}
\usepackage{xcolor}
\usepackage[colorlinks=true, urlcolor=blue!60!black, linkcolor=blue!60!black]{hyperref}
\setlength{\parindent}{0em}
\setlength{\parskip}{0.45em}
\renewcommand{\baselinestretch}{1.0}
\begin{document}
\begin{center}
{\Large \textit{United States v. Butler}, 297 U.S. 1 (1936)\par}
\vspace{0.3em}
{\normalsize Case Brief --- One-Page Edition\par}
\end{center}

\textbf{Memory Jogger.} The AAA processing-tax case: a Hamiltonian spending power, but no purchased regulation (\href{https://supreme.justia.com/cases/federal/us/297/1/}{297 U.S. 1}).

\textbf{Facts.} The Agricultural Adjustment Act paid farmers to reduce production, funded by an earmarked processing tax. The United States claimed the tax from the receivers of Hoosac Mills, a bankrupt cotton mill.

\textbf{Procedural History.} The district court ordered payment; the First Circuit reversed.

\textbf{Votes.} Affirmed, 6--3. Roberts wrote for the Court; Stone dissented, joined by Brandeis and Cardozo.

\textbf{Holding.} The general welfare power is not confined to the other enumerated powers, but Congress may not use taxing and spending to coerce or purchase regulation of local agriculture reserved to the states.

\textbf{Reasoning.} Roberts adopted Hamilton's reading but found the program coercive and a purchase of compliance Congress could not compel. Stone replied that courts judge power, not wisdom, that the program offered hope of gain rather than threat of loss, and that spending without conditions is meaningless.

\textbf{Examples.} \emph{Same-Side:} conditioning all Medicaid funds on a farm rule. \emph{Opposite-Side:} a new, optional grant. \emph{Fence-Sitter:} withholding a quarter of highway funds.

\textbf{Critique.} Critics deride the opinion's formalism and its incoherent half-Hamiltonianism.

\textbf{Key Quotation.} ``One is that courts are concerned only with the power to enact statutes, not with their wisdom. The other is that, while unconstitutional exercise of power by the executive and legislative branches of the government is subject to judicial restraint, the only check upon our own exercise of power is our own sense of self-restraint.'' --- Justice Stone, dissenting, 297 U.S.\ at 78--79.

\end{document}
